\documentclass[11pt]{article}
\usepackage[T1]{fontenc}
\usepackage[utf8]{inputenc}
\usepackage{lmodern}
\usepackage{amsmath,amssymb,amsthm}
\usepackage[margin=25mm]{geometry}
\PassOptionsToPackage{hyphens}{url}
\usepackage[hidelinks]{hyperref}
\hypersetup{pdftitle={Stieltjes moment classification of numerical Hoggatt row sums},pdfauthor={Aurantium}}
\newtheorem{theorem}{Theorem}
\newtheorem{lemma}{Lemma}
\newcommand{\Z}{\mathbb Z}
\newcommand{\Q}{\mathbb Q}
\title{Stieltjes moment classification\protect\\ of numerical Hoggatt row sums}
\author{Aurantium}
\date{14 October 2026}
\begin{document}
\maketitle

\begin{abstract}
For each positive integer $d$, we consider the numerical row sums of the
$d$-Hoggatt triangle, indexed from zero. We give a complete classification:
these row sums form a Stieltjes moment sequence exactly for $d=1,2$.
The positive cases have explicit measures. For $d=3,4$ we give finite
negative shifted-Hankel determinants. For every $d\geq5$, one degree-20
determinant identity provides an obstruction; the identity is certified
by a degree bound and 21 exact evaluations. The proof and the exact
finite certificates are supplied in full. The Baxter indexing and the
scope of earlier nonmoment assertions are stated separately.
\end{abstract}

\section{Definitions and statement}

For integers $d\geq1$, $n\geq0$ and $0\leq k\leq n$, define
\begin{equation}\label{eq:H}
 H_d(n,k)=\prod_{j=0}^{d-1}
       \frac{\binom{n+j}{k}}{\binom{k+j}{k}},
 \qquad S_d(n)=\sum_{k=0}^{n}H_d(n,k).
\end{equation}
Set $H_d(n,k)=0$ for integers $k<0$ or $k>n$.
All products in \eqref{eq:H} have positive denominators.
This is the usual normalization of the Hoggatt triangle; its entries
also equal the consecutive Pascal minors
$\det[\binom{n+i}{k+j}]_{i,j=0}^{d-1}$
\cite[Definition 1.1 and (2.1)]{MS}. The product in \eqref{eq:H}
is the definition used in every calculation below. In particular,
$S_d(0)=1$.

A sequence $(s_n)_{n\geq0}$ is a \emph{Stieltjes moment sequence}
if there is a nonnegative Borel measure $\mu$ on $[0,\infty)$ such that
\begin{equation}\label{eq:moment}
 s_n=\int_{[0,\infty)}x^n\,d\mu(x)<\infty
 \qquad(n\geq0).
\end{equation}
The zero-indexed convention is essential; the total mass is $s_0$.

\begin{theorem}\label{thm:classification}
For every integer $d\geq1$, the sequence $(S_d(n))_{n\geq0}$ is
a Stieltjes moment sequence if and only if $d=1$ or $d=2$.
\end{theorem}

The theorem concerns numerical row sums, not a coefficientwise
moment property of row polynomials. Mao and Shi pose this numerical
question after their equation (6.2) \cite[Section 6]{MS}.
The following proof uses \eqref{eq:H}, explicit measures and exact
finite certificates. It does not use a conjectured positivity property.

\section{A necessary condition and the positive cases}

For $m\geq1$, put
\begin{equation}\label{eq:shifted-Hankel}
 M_m^{(d)}=[S_d(i+j+1)]_{i,j=0}^{m-1}.
\end{equation}
If \eqref{eq:moment} holds for $S_d$, then for every real vector
$v=(v_0,\ldots,v_{m-1})^{\mathsf T}$,
\begin{equation}\label{eq:Gram}
 v^{\mathsf T}M_m^{(d)}v
 =\int_{[0,\infty)}x
       \left(\sum_{i=0}^{m-1}v_i x^i\right)^2d\mu(x)\geq0.
\end{equation}
Only finite sums are expanded. The moments up to degree $2m-1$
are finite, so each expanded term is absolutely integrable.
Thus $M_m^{(d)}$ is real symmetric positive semidefinite, and
\begin{equation}\label{eq:necessary-det}
 \det M_m^{(d)}\geq0.
\end{equation}
A negative determinant suffices to rule out a Stieltjes representation;
no converse to this necessary condition is needed.

For $d=1$, equation \eqref{eq:H} gives $H_1(n,k)=\binom nk$, hence
$S_1(n)=2^n$. The unit point mass $\delta_2$ represents these moments.

For $d=2$, factorial cancellation gives
\[
 H_2(n,k)=\frac{1}{k+1}\binom nk\binom{n+1}{k}
 =\frac{1}{n+1}\binom{n+1}{k+1}\binom{n+1}{k}.
\]
By the finite Vandermonde convolution,
\begin{equation}\label{eq:Catalan-shift}
 S_2(n)=\frac{1}{n+1}\binom{2n+2}{n}
       =\frac{1}{n+2}\binom{2n+2}{n+1}=C_{n+1}.
\end{equation}
Here $C_j=(j+1)^{-1}\binom{2j}{j}$ denotes the Catalan number.
Define the measure
\begin{equation}\label{eq:positive-measure}
 d\mu_2(x)=\frac{\sqrt{x(4-x)}}{2\pi}
               \mathbf1_{(0,4)}(x)\,dx.
\end{equation}
Its endpoint singularity exponents are both $1/2$, so it is finite;
its compact support makes every moment finite. With $x=4u$,
\begin{align}
 \int x^n\,d\mu_2(x)
 &=\frac{4^{n+2}}{2\pi}
          B\!\left(n+\frac32,\frac32\right)\notag\\
 &=\frac{(2n+2)!}{(n+1)!(n+2)!}=C_{n+1}.\label{eq:beta}
\end{align}
The beta arguments are positive. The second equality follows from
$B(a,b)=\Gamma(a)\Gamma(b)/\Gamma(a+b)$,
$\Gamma(3/2)=\sqrt\pi/2$ and
$\Gamma(n+3/2)=(2n+2)!\sqrt\pi/(4^{n+1}(n+1)!)$.
At $n=0$, the value is $1$, as required. This is the shifted Catalan
measure, with the factor $x$ relative to the usual Catalan density.

\section{Polynomial dependence on the parameter}

\begin{lemma}\label{lem:rectangle}
For nonnegative integers $d,n,k$ with $0\leq k\leq n$,
\begin{equation}\label{eq:rectangle}
 \prod_{r=0}^{d-1}\frac{\binom{n+r}{k}}{\binom{k+r}{k}}
 =\prod_{i=1}^{k}\prod_{j=1}^{n-k}
                    \frac{d+i+j-1}{i+j-1}.
\end{equation}
The left side at $d=0$ is an empty product. Empty rectangular
products at $k=0$ or $k=n$ also equal $1$.
\end{lemma}

\begin{proof}
Write $a=k$, $b=n-k$ and denote the two sides by $F_d$ and $G_d$.
Factorial cancellation gives
\[
 F_d=\prod_{r=0}^{d-1}\frac{(a+b+r)!r!}{(a+r)!(b+r)!}.
\]
Both $F_0$ and $G_0$ are $1$. For $d\geq0$,
\begin{align}
 \frac{F_{d+1}}{F_d}
 &=\frac{(a+b+d)!d!}{(a+d)!(b+d)!}
  =\prod_{i=1}^{a}\frac{d+i+b}{d+i}\notag\\
 &=\prod_{i=1}^{a}\prod_{j=1}^{b}
                    \frac{d+i+j}{d+i+j-1}
  =\frac{G_{d+1}}{G_d}.\label{eq:bridge-ratio}
\end{align}
The inner product telescopes, and every denominator is positive.
The equalities also hold when $a=0$ or $b=0$. Induction in $d$
proves \eqref{eq:rectangle}.
\end{proof}

For an indeterminate $z$, use the right side of \eqref{eq:rectangle}
to define $H_z(n,k)\in\Q[z]$. Its degree is $k(n-k)$, with
positive leading coefficient, including degree zero at an empty
rectangle. Therefore
\begin{equation}\label{eq:row-degree}
 \deg S_z(n)\leq\left\lfloor\frac{n^2}{4}\right\rfloor.
\end{equation}
The value $d=0$ is used only for this polynomial extension and its
certificate; it is not included in Theorem~\ref{thm:classification}.

Let
\begin{equation}\label{eq:D4}
 D_4(z)=\det[S_z(i+j+1)]_{i,j=0}^{3}.
\end{equation}
The entrywise degree bounds from \eqref{eq:row-degree} form the matrix
\begin{equation}\label{eq:degree-matrix}
 \begin{pmatrix}
 0&1&2&4\\1&2&4&6\\2&4&6&9\\4&6&9&12
 \end{pmatrix}.
\end{equation}
In the Leibniz expansion, the term indexed by $\pi\in\mathfrak S_4$
has degree at most
\[
 w(\pi)=\sum_{i=0}^{3}
           \left\lfloor\frac{(i+\pi(i)+1)^2}{4}\right\rfloor.
\]
All 24 bounds are listed in Appendix~\ref{app:degrees}; their maximum
is $20$. Thus $\deg D_4\leq20$, regardless of possible cancellation.

\section{The negative certificates}

\subsection{Parameters three and four}

The needed moment prefixes, evaluated by the finite product
\eqref{eq:H}, are shown below. A dash means that the value is not
needed for the displayed obstruction.
\begin{center}\small
\begin{tabular}{r|r|r}
$n$&$S_3(n)$&$S_4(n)$\\\hline
0&1&1\\
1&2&2\\
2&6&7\\
3&22&32\\
4&92&177\\
5&422&1122\\
6&2074&7898\\
7&10754&60398\\
8&58202&494078\\
9&326240&4274228\\
10&1882960&--\\
11&11140560&--
\end{tabular}
\end{center}
With the indexing of \eqref{eq:shifted-Hankel}, these prefixes give
\begin{equation}\label{eq:small-negative}
 \det M_6^{(3)}=-115894592,
 \qquad \det M_5^{(4)}=-2220868988.
\end{equation}
Each matrix is specified completely by the table and
\eqref{eq:shifted-Hankel}. The exact certificate also stores both
matrices and their evaluations by the signed-permutation formula
and rational Gaussian elimination. Their negativity contradicts
\eqref{eq:necessary-det} for $d=3,4$.

\subsection{One identity for every parameter at least five}

Define
\begin{align}
 P(z)={}&z^9+18z^8+83z^7-214z^6-2605z^5-5602z^4\notag\\
       &-1235z^3-3490z^2-10068z-1080,\label{eq:P}\\
 Q(z)={}&-\frac{z^4(z-1)^3(z+1)^3(z+2)P(z)}{74649600}.
                                                    \label{eq:Q}
\end{align}
The polynomial $Q$ has degree $4+3+3+1+9=20$.
The following exact evaluation certificate records the common value
of $D_4(d)$ and $Q(d)$ at each of the 21 distinct nodes $d=0,\ldots,20$.
\begin{center}\small
\begin{tabular}{r|r@{\qquad}r|r}
$d$&$D_4(d)=Q(d)$&$d$&$D_4(d)=Q(d)$\\\hline
0&0&11&-30625082179120\\
1&0&12&-164521620189927\\
2&5&13&-770983662611344\\
3&2832&14&-3220116854652880\\
4&163385&15&-12187081083520000\\
5&-154000&16&-42347156578932480\\
6&-152388040&17&-136534413420057600\\
7&-4142474112&18&-412026373330899147\\
8&-61044214896&19&-1172223396095416560\\
9&-620565390720&20&-3163289873712112375\\
10&-4831832968875&&
\end{tabular}
\end{center}
At each node, the exact data contain the eight row sums from index
$0$ through $7$, the two determinant evaluations, and $Q(d)$.
The supplied Python program reproduces them from \eqref{eq:H}.
Since $D_4-Q$ has degree at most $20$ and has these 21 distinct zeros,
\begin{equation}\label{eq:identity}
 D_4(z)=Q(z)\qquad\hbox{in }\Q[z].
\end{equation}
This is a polynomial identity certificate, not extrapolation of a
finite collection of signs. No parameter-dependent factor is divided
out at the nodes $0$ or $1$.

The exact binomial expansion at $z=t+5$ is
\begin{align}
 P(t+5)={}&t^9+63t^8+1703t^7+25791t^6+239300t^5\notag\\
 &+1393398t^4+4978600t^3+10009560t^2\notag\\
 &+8752032t+190080.\label{eq:positive-shift}
\end{align}
Every coefficient is positive, including the constant term.
Thus $P(d)>0$ whenever $d\geq5$. For such $d$, all remaining
factors of $-Q(d)$ in \eqref{eq:Q} are positive. Consequently
\begin{equation}\label{eq:all-negative}
 \det M_4^{(d)}=D_4(d)<0\qquad(d\in\Z,\ d\geq5).
\end{equation}

For a concrete member of this uniform obstruction,
\[
 M_4^{(5)}=\begin{pmatrix}
 2&8&44&310\\
 8&44&310&2606\\
 44&310&2606&25202\\
 310&2606&25202&272582
 \end{pmatrix},\qquad \det M_4^{(5)}=-154000.
\]
The redundant integer-vector certificate
\begin{equation}\label{eq:quadratic}
 \begin{aligned}
 v&=(-17017,11284,-1877,85)^{\mathsf T},\\
 M_4^{(5)}v&=(0,0,0,-3850)^{\mathsf T},\\
 v^{\mathsf T}M_4^{(5)}v&=-327250
 \end{aligned}
\end{equation}
gives the same obstruction directly through \eqref{eq:Gram}.

The measures for $d=1,2$, the two determinants in
\eqref{eq:small-negative}, and \eqref{eq:all-negative} cover every
integer $d\geq1$. This proves Theorem~\ref{thm:classification}.

\section{Baxter indexing and relation to earlier work}

Let $\beta_n$, $n\geq1$, denote the Baxter numbers in equation (1.1)
of Fang, Zhang and Zhao \cite{FZZ}, with first terms
$\beta_1=1,\beta_2=2,\beta_3=6,\beta_4=22$.
Their summand satisfies
\[
 D_{n+1,k+1}
 =\frac{2n!(n+1)!(n+2)!}
 {k!(k+1)!(k+2)!(n-k)!(n+1-k)!(n+2-k)!}
 =H_3(n,k).
\]
Hence the present convention is $S_3(n)=\beta_{n+1}$.
The shift was already made explicit by Fielder and Alford
\cite[p.~162]{FA}. Mao and Shi instead write $B_n=S_3(n)$ and
specify $B_0=1,B_1=2$ \cite[pp.~16--17]{MS}.

Fang, Zhang and Zhao explicitly state that their positively indexed
Baxter sequence is not Stieltjes, immediately after Conjecture 4.2
of their preprint \cite[p.~11]{FZZ}. This prior assertion must be
distinguished from their accompanying reference to a negative
fifth-order leading Hankel minor: its precise starting index and
alignment with the zero-indexed sequence here have not been
unambiguously established. Their final journal text has not been
checked here, and no correction by those authors is asserted.
We give the explicit normalized matrices above without claiming
first discovery of Baxter non-Stieltjesness.

This distinction matters because if $(a_n)_{n\geq0}$ has measure
$\mu$, its forward shift has measure $x\,d\mu$; non-Stieltjesness
does not propagate forward by that implication. A negative certificate
for a sequence with an additionally chosen initial term cannot simply
be transferred after that term is deleted. For general definitions
of moment sequences and their Hankel matrices, see \cite{WZ}.

The strong coefficientwise $q$-log-convexity result of Chern and Shi
\cite[Theorem 1.2]{CS} concerns a different positivity property and
does not supply a Stieltjes representation for these numerical
row sums. The literature priority of the full classification and
of the particular uniform determinant formula is not certified here.

\section{Exact data and reproduction}

The file \texttt{data/exact\_certificate.json} contains the 24
permutation degree bounds, the 21 evaluation records, all three
negative matrices, the vector certificate, and the ten coefficients
of \eqref{eq:positive-shift}. Its arithmetic is integer and rational;
no floating-point determinant or numerical sign threshold is used.
The program \texttt{code/verify\_certificate.py} reconstructs the
finite values from \eqref{eq:H}, cross-checks the rectangular product,
and compares the signed-permutation determinant with rational
Gaussian elimination. The argument establishing \eqref{eq:identity}
also requires the degree proof above, and the positive cases require
their measure calculation; finite tests alone do not replace either.

The saved certificate comes from an earlier successful execution of
the underlying arithmetic. The stdout-only packaged entry also completed one local replay with
Python 3.12.10 on 14 October 2026. All checks passed, its JSON output
matched the archived certificate, and child, guard and shell exited
zero with empty stderr. No original
author full suite is claimed to have been rerun. The precise historical
record and the minimal command are described in
\texttt{REPRODUCIBILITY.md}.

\paragraph{Acknowledgement.}
This note was prepared with AI assistance.

\appendix
\section{Complete degree-bound certificate}\label{app:degrees}
The following table lists every permutation of $\{0,1,2,3\}$ and
the bound $w(\pi)$ used with \eqref{eq:degree-matrix}.
\begin{center}\small
\begin{tabular}{c|r@{\qquad}c|r}
$\pi$&$w(\pi)$&$\pi$&$w(\pi)$\\\hline
$(0,1,2,3)$&20&$(2,0,1,3)$&19\\
$(0,1,3,2)$&20&$(2,0,3,1)$&18\\
$(0,2,1,3)$&20&$(2,1,0,3)$&18\\
$(0,2,3,1)$&19&$(2,1,3,0)$&17\\
$(0,3,1,2)$&19&$(2,3,0,1)$&16\\
$(0,3,2,1)$&18&$(2,3,1,0)$&16\\
$(1,0,2,3)$&20&$(3,0,1,2)$&18\\
$(1,0,3,2)$&20&$(3,0,2,1)$&17\\
$(1,2,0,3)$&19&$(3,1,0,2)$&17\\
$(1,2,3,0)$&18&$(3,1,2,0)$&16\\
$(1,3,0,2)$&18&$(3,2,0,1)$&16\\
$(1,3,2,0)$&17&$(3,2,1,0)$&16
\end{tabular}
\end{center}

\begin{thebibliography}{9}\footnotesize
\setlength{\itemsep}{2pt}
\raggedright
\bibitem{FA}
D.~C.~Fielder and C.~O.~Alford,
On a conjecture by Hoggatt with extensions to Hoggatt sums and
Hoggatt triangles, \emph{Fibonacci Quart.} \textbf{27}(2)
(1989), 160--168.
\url{https://www.mathstat.dal.ca/FQ/Scanned/27-2/fielder.pdf}.

\bibitem{MS}
J.~Mao and W.~Shi, \emph{The analytic properties of Hoggatt triangles},
arXiv:2607.01582v1, 2 July 2026. Definition 1.1, (2.1), and Section 6.
\url{https://arxiv.org/abs/2607.01582v1}.

\bibitem{FZZ}
H.~Fang, C.~X.~T.~Zhang and J.~J.~Y.~Zhao,
\emph{Analytic properties arising from the Baxter numbers},
arXiv:2505.05873v1, submitted 9 May 2025, Section 4, p.~11.
\url{https://arxiv.org/abs/2505.05873v1}.
Journal record: \emph{J. Math. Anal. Appl.} \textbf{561}(1)
(2026), 130615, \url{https://doi.org/10.1016/j.jmaa.2026.130615}.
The discussion here is tied to the specified preprint; the final
journal text has not been checked.

\bibitem{WZ}
Y.~Wang and B.-X.~Zhu, Log-convex and Stieltjes moment sequences,
\emph{Adv. in Appl. Math.} \textbf{81} (2016), 115--127.
\url{https://arxiv.org/abs/1612.04114v1}.

\bibitem{CS}
S.~Chern and W.~Shi, \emph{Strong $q$-log-convexity of the
$d$-Hoggatt polynomials}, arXiv:2609.33732v1, 27 September 2026,
Theorem 1.2.
\url{https://arxiv.org/abs/2609.33732v1}.
\end{thebibliography}
\end{document}
